\section{Introdução}
\label{cap:introducao}

A Teoria da Probabilidade é um dos principais ramos da Matemática Discreta voltados ao
estudo de eventos aleatórios em espaços amostrais finitos ou contáveis. Seu objetivo é
quantificar a chance de ocorrência de determinados eventos por meio de modelos
matemáticos, permitindo que situações incertas sejam analisadas de forma lógica e objetiva.
Essa teoria possui aplicações em diversas áreas, como computação, inteligência artificial,
ciência de dados, economia, engenharia e segurança da informação, sendo um importante
instrumento para a tomada de decisões fundamentadas.

Na Matemática Discreta, a probabilidade está diretamente relacionada à análise combinatória
e aos conjuntos, pois muitos problemas envolvem a contagem de possibilidades para
determinar a probabilidade de um evento ocorrer. A abordagem clássica define a
probabilidade como a razão entre o número de casos favoráveis ($m$) e o número total de
casos possíveis ($n$), desde que todos possuam a mesma chance de ocorrência:

\[ P(A)=\frac{m}{n} \]

Essa definição constitui a base para o desenvolvimento de técnicas probabilísticas mais avançadas.

% TODO - adicionar conteúdo aqui. Essa introdução está chucra
% TODO - A introdução precisa informar brevemente o objetivo, a metodologia e os principais resultados a serem apresentados.

% TODO - Na Introdução deve-se apresentar a proposta do trabalho com referencial teórico. Termine a seção descrevendo a organização/estrutura do trabalho/relatório. Nenhum parágrafo deve ter menos de 3 linhas e a seção deve ter pelo menos uma página inteira (incluindo figuras e tabelas se houver) ou pelo menos 4 parágrafos.

% TODO - o texto abaixo faz algum sentido?
Este trabalho tem como objetivo apresentar os fundamentos da Teoria da Probabilidade e
sua relação com esses conceitos complementares, evidenciando sua importância em
aplicações computacionais e na resolução de problemas discretos. Para isso, o relatório está
organizado da seguinte forma: inicialmente são descritos os materiais e a metodologia
utilizada; em seguida, são apresentados os resultados obtidos e suas respectivas análises;
por fim, são apresentadas as conclusões decorrentes do estudo realizado. (???)
